% fig/pipeline.tex — 時空間イベント帰属パイプライン図
% % fig/pipeline.tex — 時空間イベント帰属パイプライン図
% % fig/pipeline.tex — 時空間イベント帰属パイプライン図
% % fig/pipeline.tex — 時空間イベント帰属パイプライン図
% \input{fig/pipeline} で読み込む
% 前提: プリアンブルで \usepackage{tikz} および
%       \usetikzlibrary{arrows.meta, positioning} を宣言済み

\begin{figure}[t]
\centering
\begin{tikzpicture}[
  node distance=0.4cm,
  box/.style={
    draw, rounded corners=2pt, thick,
    minimum width=5.6cm, minimum height=0.85cm,
    text width=5.4cm, align=center,
    font=\small
  },
  arr/.style={-{Stealth[length=2.5mm]}, thick},
]

% Step 1
\node[box] (s1) {%
  \textbf{Step 1:} サポート曲面の計算\\[-1pt]
  {\scriptsize トランザクション＋位置 $\to$ $S_P(t,v)$}%
};

% Step 2
\node[box, below=of s1] (s2) {%
  \textbf{Step 2:} 位置別変化点検出\\[-1pt]
  {\scriptsize 各位置で閾値交差を検出}%
};

% Step 3
\node[box, below=of s2] (s3) {%
  \textbf{Step 3:} 時空間帰属スコアリング\\[-1pt]
  {\scriptsize $\mathrm{prox}_t \times \mathrm{prox}_s \times \mathrm{mag}$}%
};

% Step 4
\node[box, below=of s3] (s4) {%
  \textbf{Step 4:} 時空間置換検定\\[-1pt]
  {\scriptsize 循環時間シフト＋空間範囲保持}%
};

% Step 5
\node[box, below=of s4] (s5) {%
  \textbf{Step 5:} グローバルBH補正\\[-1pt]
  {\scriptsize FDR制御}%
};

% Step 6
\node[box, below=of s5] (s6) {%
  \textbf{Step 6:} アイテム重複排除\\[-1pt]
  {\scriptsize Union-Find}%
};

% Arrows
\draw[arr] (s1) -- (s2);
\draw[arr] (s2) -- (s3);
\draw[arr] (s3) -- (s4);
\draw[arr] (s4) -- (s5);
\draw[arr] (s5) -- (s6);

\end{tikzpicture}
\caption{提案する時空間イベント帰属パイプラインの全体構成。
6段のステップにより、サポート曲面の計算から
有意な帰属結果の出力までを一貫して処理する。}
\label{fig:pipeline}
\end{figure}
 で読み込む
% 前提: プリアンブルで \usepackage{tikz} および
%       \usetikzlibrary{arrows.meta, positioning} を宣言済み

\begin{figure}[t]
\centering
\begin{tikzpicture}[
  node distance=0.4cm,
  box/.style={
    draw, rounded corners=2pt, thick,
    minimum width=5.6cm, minimum height=0.85cm,
    text width=5.4cm, align=center,
    font=\small
  },
  arr/.style={-{Stealth[length=2.5mm]}, thick},
]

% Step 1
\node[box] (s1) {%
  \textbf{Step 1:} サポート曲面の計算\\[-1pt]
  {\scriptsize トランザクション＋位置 $\to$ $S_P(t,v)$}%
};

% Step 2
\node[box, below=of s1] (s2) {%
  \textbf{Step 2:} 位置別変化点検出\\[-1pt]
  {\scriptsize 各位置で閾値交差を検出}%
};

% Step 3
\node[box, below=of s2] (s3) {%
  \textbf{Step 3:} 時空間帰属スコアリング\\[-1pt]
  {\scriptsize $\mathrm{prox}_t \times \mathrm{prox}_s \times \mathrm{mag}$}%
};

% Step 4
\node[box, below=of s3] (s4) {%
  \textbf{Step 4:} 時空間置換検定\\[-1pt]
  {\scriptsize 循環時間シフト＋空間範囲保持}%
};

% Step 5
\node[box, below=of s4] (s5) {%
  \textbf{Step 5:} グローバルBH補正\\[-1pt]
  {\scriptsize FDR制御}%
};

% Step 6
\node[box, below=of s5] (s6) {%
  \textbf{Step 6:} アイテム重複排除\\[-1pt]
  {\scriptsize Union-Find}%
};

% Arrows
\draw[arr] (s1) -- (s2);
\draw[arr] (s2) -- (s3);
\draw[arr] (s3) -- (s4);
\draw[arr] (s4) -- (s5);
\draw[arr] (s5) -- (s6);

\end{tikzpicture}
\caption{提案する時空間イベント帰属パイプラインの全体構成。
6段のステップにより、サポート曲面の計算から
有意な帰属結果の出力までを一貫して処理する。}
\label{fig:pipeline}
\end{figure}
 で読み込む
% 前提: プリアンブルで \usepackage{tikz} および
%       \usetikzlibrary{arrows.meta, positioning} を宣言済み

\begin{figure}[t]
\centering
\begin{tikzpicture}[
  node distance=0.4cm,
  box/.style={
    draw, rounded corners=2pt, thick,
    minimum width=5.6cm, minimum height=0.85cm,
    text width=5.4cm, align=center,
    font=\small
  },
  arr/.style={-{Stealth[length=2.5mm]}, thick},
]

% Step 1
\node[box] (s1) {%
  \textbf{Step 1:} サポート曲面の計算\\[-1pt]
  {\scriptsize トランザクション＋位置 $\to$ $S_P(t,v)$}%
};

% Step 2
\node[box, below=of s1] (s2) {%
  \textbf{Step 2:} 位置別変化点検出\\[-1pt]
  {\scriptsize 各位置で閾値交差を検出}%
};

% Step 3
\node[box, below=of s2] (s3) {%
  \textbf{Step 3:} 時空間帰属スコアリング\\[-1pt]
  {\scriptsize $\mathrm{prox}_t \times \mathrm{prox}_s \times \mathrm{mag}$}%
};

% Step 4
\node[box, below=of s3] (s4) {%
  \textbf{Step 4:} 時空間置換検定\\[-1pt]
  {\scriptsize 循環時間シフト＋空間範囲保持}%
};

% Step 5
\node[box, below=of s4] (s5) {%
  \textbf{Step 5:} グローバルBH補正\\[-1pt]
  {\scriptsize FDR制御}%
};

% Step 6
\node[box, below=of s5] (s6) {%
  \textbf{Step 6:} アイテム重複排除\\[-1pt]
  {\scriptsize Union-Find}%
};

% Arrows
\draw[arr] (s1) -- (s2);
\draw[arr] (s2) -- (s3);
\draw[arr] (s3) -- (s4);
\draw[arr] (s4) -- (s5);
\draw[arr] (s5) -- (s6);

\end{tikzpicture}
\caption{提案する時空間イベント帰属パイプラインの全体構成。
6段のステップにより、サポート曲面の計算から
有意な帰属結果の出力までを一貫して処理する。}
\label{fig:pipeline}
\end{figure}
 で読み込む
% 前提: プリアンブルで \usepackage{tikz} および
%       \usetikzlibrary{arrows.meta, positioning} を宣言済み

\begin{figure}[t]
\centering
\begin{tikzpicture}[
  node distance=0.4cm,
  box/.style={
    draw, rounded corners=2pt, thick,
    minimum width=5.6cm, minimum height=0.85cm,
    text width=5.4cm, align=center,
    font=\small
  },
  arr/.style={-{Stealth[length=2.5mm]}, thick},
]

% Step 1
\node[box] (s1) {%
  \textbf{Step 1:} サポート曲面の計算\\[-1pt]
  {\scriptsize トランザクション＋位置 $\to$ $S_P(t,v)$}%
};

% Step 2
\node[box, below=of s1] (s2) {%
  \textbf{Step 2:} 位置別変化点検出\\[-1pt]
  {\scriptsize 各位置で閾値交差を検出}%
};

% Step 3
\node[box, below=of s2] (s3) {%
  \textbf{Step 3:} 時空間帰属スコアリング\\[-1pt]
  {\scriptsize $\mathrm{prox}_t \times \mathrm{prox}_s \times \mathrm{mag}$}%
};

% Step 4
\node[box, below=of s3] (s4) {%
  \textbf{Step 4:} 時空間置換検定\\[-1pt]
  {\scriptsize 循環時間シフト＋空間範囲保持}%
};

% Step 5
\node[box, below=of s4] (s5) {%
  \textbf{Step 5:} グローバルBH補正\\[-1pt]
  {\scriptsize FDR制御}%
};

% Step 6
\node[box, below=of s5] (s6) {%
  \textbf{Step 6:} アイテム重複排除\\[-1pt]
  {\scriptsize Union-Find}%
};

% Arrows
\draw[arr] (s1) -- (s2);
\draw[arr] (s2) -- (s3);
\draw[arr] (s3) -- (s4);
\draw[arr] (s4) -- (s5);
\draw[arr] (s5) -- (s6);

\end{tikzpicture}
\caption{提案する時空間イベント帰属パイプラインの全体構成。
6段のステップにより、サポート曲面の計算から
有意な帰属結果の出力までを一貫して処理する。}
\label{fig:pipeline}
\end{figure}
